\documentclass[12pt,a4paper]{article}
\usepackage[utf8]{inputenc}
\usepackage[T1]{fontenc}
\usepackage[french]{babel}
\usepackage{geometry}
\usepackage{graphicx}
\usepackage{booktabs}
\usepackage{longtable}
\usepackage{array}
\usepackage{hyperref}
\usepackage{listings}
\usepackage{xcolor}
\usepackage{enumitem}
\usepackage{fancyhdr}
\usepackage{titlesec}

\geometry{margin=2.5cm}
\hypersetup{colorlinks=true, linkcolor=blue!60!black, urlcolor=blue!60!black}

\definecolor{codeblue}{RGB}{0,102,204}
\definecolor{codegray}{RGB}{128,128,128}
\definecolor{codegreen}{RGB}{0,128,0}
\definecolor{backcolour}{RGB}{245,245,245}

\lstset{
    backgroundcolor=\color{backcolour},
    basicstyle=\ttfamily\small,
    breaklines=true,
    frame=single,
    rulecolor=\color{codegray},
    commentstyle=\color{codegreen},
    keywordstyle=\color{codeblue},
}

\pagestyle{fancy}
\fancyhf{}
\rhead{ESB-Learning}
\lhead{Architecture Agentic AI}
\rfoot{\thepage}

\title{\textbf{Architecture Agentic AI}\\[0.5em]
\large Plateforme ESB-Learning\\[0.3em]
\normalsize Description complète des agents, outils et pipelines}
\author{Aymen Benbrik}
\date{\today}

\begin{document}
\maketitle
\tableofcontents
\newpage

% ============================================================================
\section{Vue d'ensemble}

La plateforme ESB-Learning implémente une architecture \textbf{Agentic AI multi-agents} orchestrée par des modèles de langage (LLM). Le système comporte \textbf{9 agents/services AI} interconnectés, utilisant principalement \textbf{Google Gemini} comme LLM d'orchestration et \textbf{LangGraph} comme framework d'agents.

\subsection{Technologies clés}

\begin{table}[h!]
\centering
\begin{tabular}{@{}ll@{}}
\toprule
\textbf{Technologie} & \textbf{Rôle} \\
\midrule
Google Gemini 2.5-Flash & LLM principal (orchestration, génération, analyse) \\
Google Gemini 2.5-Pro & LLM avancé (examens, analyse approfondie) \\
LangGraph & Framework d'agents (StateGraph, ReAct) \\
TunBERT (tunis-ai/TunBERT) & Modèle BERT pré-entraîné sur le dialecte tunisien \\
ChromaDB & Base vectorielle pour RAG \\
gTTS & Synthèse vocale (Text-to-Speech) \\
PyTorch + Transformers & Inférence TunBERT \\
\bottomrule
\end{tabular}
\caption{Stack technologique IA}
\end{table}

% ============================================================================
\section{Agent Conversationnel (Assistant Agent)}

\begin{itemize}[leftmargin=*]
    \item \textbf{Fichier} : \texttt{app/services/assistant\_agent.py}
    \item \textbf{API} : \texttt{POST /api/v1/assistant/chat}
    \item \textbf{Pattern} : LangGraph ReAct Agent (function calling)
    \item \textbf{LLM} : Gemini 2.5-Flash (temperature: 0.4)
\end{itemize}

\subsection{Description}
Agent conversationnel intelligent qui s'adapte au rôle de l'utilisateur (étudiant, enseignant, administrateur) et communique en \textbf{français}, \textbf{anglais} et \textbf{dialecte tunisien}. Il utilise le pattern ReAct (Reasoning + Acting) pour décider quels outils appeler en fonction de la requête.

\subsection{Outils communs (tous les rôles)}

\begin{longtable}{@{}p{3.8cm}p{6cm}p{2cm}p{3cm}@{}}
\toprule
\textbf{Outil} & \textbf{Description} & \textbf{Entrée} & \textbf{Sortie} \\
\midrule
\endhead
\texttt{get\_my\_courses} & Récupère les cours de l'utilisateur & user\_id & Liste JSON des cours \\
\texttt{get\_calendar\_activities} & Calendrier des activités à venir & user\_id & Activités avec dates \\
\texttt{get\_course\_details} & Informations détaillées d'un cours & course\_id & Métadonnées du cours \\
\bottomrule
\end{longtable}

\subsection{Outils étudiants}

\begin{longtable}{@{}p{4cm}p{6cm}p{2cm}p{3cm}@{}}
\toprule
\textbf{Outil} & \textbf{Description} & \textbf{Entrée} & \textbf{Sortie} \\
\midrule
\endhead
\texttt{get\_my\_performance} & Scores quiz, moyennes, Bloom & student\_id & Performance par cours \\
\texttt{get\_my\_grades\_summary} & Notes par module & student\_id & Résumé des notes \\
\texttt{get\_recommendations} & Recommandations IA (Coach) & student\_id & Lacunes et plan d'étude \\
\bottomrule
\end{longtable}

\subsection{Outils enseignants}

\begin{longtable}{@{}p{4.5cm}p{5.5cm}p{2cm}p{3cm}@{}}
\toprule
\textbf{Outil} & \textbf{Description} & \textbf{Entrée} & \textbf{Sortie} \\
\midrule
\endhead
\texttt{get\_at\_risk\_students} & Détection étudiants en difficulté & teacher\_id & Liste à risque \\
\texttt{get\_class\_performance} & Statistiques globales classe & teacher\_id & Analytiques classe \\
\texttt{suggest\_quiz\_for\_student} & Quiz adapté pour un étudiant & student\_id & Quiz personnalisé \\
\bottomrule
\end{longtable}

\subsection{Flux d'exécution}

\begin{lstlisting}[language={},caption={Pipeline de l'Assistant Agent}]
Message utilisateur
  |-> Detection de langue (keyword-based)
  |-> [Si tunisien] Enrichissement TunBERT (intents semantiques)
  |-> Construction du prompt systeme (role, nom, user_id)
  |-> Selection des outils selon le role
  |-> Agent ReAct (Gemini function calling)
  |     |-> Raisonnement : analyse la requete
  |     |-> Action : appel d'un ou plusieurs outils
  |     |-> Observation : resultat de l'outil
  |     '-> Reponse finale : synthese en langage naturel
  '-> Retour : {response, language, tools_used, tunbert_intents}
\end{lstlisting}

\subsection{Endpoints associés}

\begin{itemize}
    \item \texttt{POST /api/v1/assistant/chat} --- Chat principal
    \item \texttt{POST /api/v1/assistant/tts} --- Synthèse vocale (gTTS)
    \item \texttt{POST /api/v1/assistant/stt} --- Reconnaissance vocale (Gemini audio)
    \item \texttt{GET /api/v1/assistant/tunbert-status} --- État du modèle TunBERT
\end{itemize}

% ============================================================================
\section{Service TunBERT (Dialecte tunisien)}

\begin{itemize}[leftmargin=*]
    \item \textbf{Fichier} : \texttt{app/services/tunbert\_service.py}
    \item \textbf{Pattern} : Singleton lazy-loaded + Classification par similarité cosinus
    \item \textbf{Modèle} : tunis-ai/TunBERT (BERT-base, 768 dimensions, 440 Mo)
\end{itemize}

\subsection{Description}
Service de compréhension du dialecte tunisien utilisant le modèle TunBERT pré-entraîné. Intégré comme \textbf{agent de preprocessing} dans le pipeline de l'assistant conversationnel. Quand un message en dialecte tunisien est détecté, TunBERT analyse l'intention éducative et enrichit le prompt envoyé à Gemini.

\subsection{Intentions éducatives supportées (9)}

\begin{longtable}{@{}p{4cm}p{5cm}p{5.5cm}@{}}
\toprule
\textbf{Intention} & \textbf{Description} & \textbf{Exemples tunisiens} \\
\midrule
\endhead
\texttt{ask\_courses} & Informations sur les cours & ``chnou el cours mte3i'' \\
\texttt{ask\_grades} & Notes et résultats & ``chnou el notes mte3i'' \\
\texttt{ask\_schedule} & Calendrier et horaires & ``wakteh el exam'' \\
\texttt{ask\_recommendations} & Conseils d'étude & ``chnou lazem na9ra'' \\
\texttt{ask\_performance} & Niveau et progression & ``kif niveau mte3i'' \\
\texttt{greeting} & Salutation tunisienne & ``ahla bik'', ``salam 3likom'' \\
\texttt{ask\_help} & Aide sur l'utilisation & ``3awni'', ``ki nesta3mel'' \\
\texttt{teacher\_at\_risk} & Étudiants en difficulté & ``chkoun el talaba fi danger'' \\
\texttt{teacher\_class\_perf} & Performance de la classe & ``kif el classe'' \\
\bottomrule
\end{longtable}

\subsection{Architecture technique}

\begin{lstlisting}[language={},caption={Chargement et classification TunBERT}]
Premiere requete tunisienne
  |-> Chargement lazy du modele (440 Mo, une seule fois)
  |-> Remapping des cles (strip "BertModel." prefix)
  |-> Pre-calcul des embeddings d'ancrage (9 intentions x 5-7 phrases)
  '-> Modele pret en memoire (singleton)

Requete suivante
  |-> Tokenisation du texte (BertTokenizer)
  |-> Embedding CLS (768 dimensions)
  |-> Similarite cosinus avec les ancrages
  |-> Seuil de confiance (0.75)
  '-> Enrichissement du prompt si confiance suffisante
\end{lstlisting}

% ============================================================================
\section{Coach Agent (Analyse de performance)}

\begin{itemize}[leftmargin=*]
    \item \textbf{Fichier} : \texttt{app/services/coach\_agent.py}
    \item \textbf{API} : \texttt{GET /api/v1/coach/analyze}
    \item \textbf{Pattern} : LangGraph StateGraph (pipeline 4 étapes)
    \item \textbf{LLM} : Gemini 2.5-Flash (temperature: 0.3)
\end{itemize}

\subsection{Description}
Agent d'analyse des performances étudiantes. Collecte les données de quiz et examens, identifie les lacunes par compétence et niveau de Bloom, puis génère des recommandations personnalisées et un plan d'étude.

\subsection{Pipeline}

\begin{table}[h!]
\centering
\begin{tabular}{@{}clp{8cm}@{}}
\toprule
\textbf{\#} & \textbf{Étape} & \textbf{Description} \\
\midrule
1 & Collecte des données & Agrégation des scores quiz par cours, calcul Bloom \\
2 & Analyse LLM & Envoi des données à Gemini pour analyse approfondie \\
3 & Génération recommandations & skill\_gaps, recommendations, study\_plan \\
4 & Formatage sortie & JSON structuré avec niveaux de sévérité \\
\bottomrule
\end{tabular}
\caption{Pipeline du Coach Agent}
\end{table}

\subsection{Sortie structurée}
L'agent retourne un JSON contenant :
\begin{itemize}
    \item \textbf{performance} : scores par cours, moyenne globale, scores Bloom
    \item \textbf{skill\_gaps} : lacunes identifiées avec sévérité (high/medium/low)
    \item \textbf{recommendations} : exercices recommandés avec priorité (urgent/important/optional)
    \item \textbf{study\_plan} : planning d'étude avec activités et durées
\end{itemize}

% ============================================================================
\section{Exam Agent Graph (Évaluation multi-agents)}

\begin{itemize}[leftmargin=*]
    \item \textbf{Fichier} : \texttt{app/services/exam\_agent\_graph.py}
    \item \textbf{API} : \texttt{POST /courses/<id>/exam/analyze}
    \item \textbf{Pattern} : LangGraph StateGraph (pipeline 10 étapes asynchrone)
    \item \textbf{LLM} : Gemini 2.5-Flash / Pro
\end{itemize}

\subsection{Description}
Système multi-agents pour l'analyse pédagogique complète d'un examen. Extrait les questions, les classifie par AA et taxonomie de Bloom, compare avec le syllabus, génère des feedbacks et compile un document \LaTeX.

\subsection{Pipeline (10 étapes)}

\begin{longtable}{@{}clp{4cm}p{6cm}@{}}
\toprule
\textbf{\#} & \textbf{Étape} & \textbf{Fonction} & \textbf{Description} \\
\midrule
\endhead
1 & Extraction texte & \texttt{extract\_text} & PDF/DOCX $\rightarrow$ texte brut \\
2 & Extraction questions & \texttt{extract\_questions} & Parsing avec points \\
3 & Classification AA & \texttt{classify\_aa} & Mapping Acquis d'Apprentissage \\
4 & Classification Bloom & \texttt{classify\_bloom} & Niveau taxonomique \\
5 & Difficulté & \texttt{assess\_difficulty} & Difficulté relative au cours \\
6 & Comparaison contenu & \texttt{compare\_content} & Alignement Syllabus $\leftrightarrow$ Examen \\
7 & Analyse feedback & \texttt{analyze\_feedback} & Feedback pédagogique \\
8 & Suggestions & \texttt{suggest\_adjustments} & Améliorations par question \\
9 & Génération \LaTeX & \texttt{generate\_latex} & Document + compilation PDF \\
10 & Évaluation & \texttt{evaluate\_proposal} & Évaluation pédagogique finale \\
\bottomrule
\end{longtable}

\subsection{Persistance}
\begin{itemize}
    \item Progression sauvegardée dans \texttt{ExamAnalysisSession.progress} (0-100\%)
    \item Label de l'agent courant pour l'UI
    \item Questions extraites persistées dans \texttt{ExamExtractedQuestion}
    \item État complet sauvegardé en JSON dans \texttt{state\_json}
\end{itemize}

% ============================================================================
\section{TP Agent Graph (Travaux Pratiques)}

\begin{itemize}[leftmargin=*]
    \item \textbf{Fichier} : \texttt{app/services/tp\_agent\_graph.py}
    \item \textbf{Pattern} : LangGraph StateGraph (2 workflows distincts)
    \item \textbf{LLM} : Gemini 2.5-Flash
\end{itemize}

\subsection{Workflow Création (Enseignant --- 5 étapes)}

\begin{table}[h!]
\centering
\begin{tabular}{@{}clp{8cm}@{}}
\toprule
\textbf{\#} & \textbf{Étape} & \textbf{Description} \\
\midrule
1 & \texttt{get\_context} & Récupération du contexte de la section \\
2 & \texttt{generate\_statement} & Génération IA de l'énoncé du TP \\
3 & \texttt{parse\_questions} & Extraction des questions structurées \\
4 & \texttt{suggest\_aa} & Suggestion des codes AA \\
5 & \texttt{generate\_reference} & Solution de référence + critères \\
\bottomrule
\end{tabular}
\caption{Workflow de création de TP}
\end{table}

\subsection{Workflow Correction (Automatique --- 2 étapes)}

\begin{table}[h!]
\centering
\begin{tabular}{@{}clp{8cm}@{}}
\toprule
\textbf{\#} & \textbf{Étape} & \textbf{Description} \\
\midrule
1 & \texttt{auto\_correct} & Analyse soumission vs référence \\
2 & \texttt{propose\_grade} & Finalisation de la note numérique \\
\bottomrule
\end{tabular}
\caption{Workflow de correction automatique}
\end{table}

\subsection{Outils}

\begin{longtable}{@{}p{5cm}p{9cm}@{}}
\toprule
\textbf{Outil} & \textbf{Description} \\
\midrule
\endhead
\texttt{get\_section\_context} & Récupération contenu éducatif \\
\texttt{generate\_tp\_statement} & Génération d'énoncé \\
\texttt{parse\_tp\_questions} & Extraction de questions \\
\texttt{suggest\_aa\_codes} & Suggestion AA \\
\texttt{generate\_reference\_solution} & Solution de référence \\
\texttt{auto\_correct\_submission} & Correction automatique \\
\texttt{propose\_grade} & Proposition de note \\
\texttt{chat\_with\_student} & Tutorat socratique (TP formatifs) \\
\bottomrule
\end{longtable}

\textbf{Langages supportés} : Python, SQL, R, Java, C, C++

% ============================================================================
\section{Service de Feedback (Post-Évaluation)}

\begin{itemize}[leftmargin=*]
    \item \textbf{Fichier} : \texttt{app/services/feedback\_service.py}
    \item \textbf{API} : \texttt{POST /api/v1/feedback/generate/<exam\_session\_id>}
    \item \textbf{Pattern} : Appel LLM direct (JSON output)
    \item \textbf{LLM} : Gemini 2.5-Flash (temperature: 0.4)
\end{itemize}

Génère automatiquement un feedback pédagogique personnalisé après chaque épreuve : forces, faiblesses, recommandations, et feedback formaté en Markdown.

% ============================================================================
\section{Services RAG (Chat documentaire)}

\begin{itemize}[leftmargin=*]
    \item \textbf{Fichiers} : \texttt{chat\_service.py}, \texttt{vector\_store.py}, \texttt{document\_pipeline.py}
    \item \textbf{API} : \texttt{POST /api/v1/ai/chat/<document\_id>}
    \item \textbf{Pattern} : RAG (Retrieval-Augmented Generation) avec CAG
    \item \textbf{Base vectorielle} : ChromaDB
\end{itemize}

Système de chat intelligent sur les documents de cours. Pipeline : CAG (recherche résumés) $\rightarrow$ RAG (recherche détaillée) $\rightarrow$ Génération LLM avec citations.

% ============================================================================
\section{Services d'Extraction et Traitement}

\subsection{Smart Extraction Service}
Extraction intelligente de contenu pour la génération de quiz, alignée sur les CLO et objectifs. Scoring par type de section (définitions: 2.0x, formules: 1.7x, exemples: 1.4x).

\subsection{Program Extraction Service}
Pipeline Agentic AI pour l'extraction automatique depuis le descriptif de formation : AAP, compétences, modules (par semestre), création automatique des enseignants, affectation aux cours.

\subsection{Video Analysis Service}
Analyse vidéo avec Gemini Vision : extraction de frames, transcription audio, analyse visuelle, timeline, rapport PDF.

\subsection{Syllabus Service}
Extraction des CLO, plan hebdomadaire, PLO, TNE depuis les PDF syllabus.

% ============================================================================
\section{Tableau récapitulatif des agents}

\begin{longtable}{@{}p{2.5cm}p{2cm}p{2cm}p{2.5cm}cp{3cm}@{}}
\toprule
\textbf{Agent} & \textbf{Type} & \textbf{Modèle} & \textbf{Pattern} & \textbf{Étapes} & \textbf{Point d'entrée} \\
\midrule
\endhead
Assistant & Conversationnel & Gemini Flash & ReAct & 1 & \texttt{/assistant/chat} \\
TunBERT & Preprocessing & TunBERT & Cosine sim. & 1 & Appelé par Assistant \\
Coach & Analytique & Gemini Flash & StateGraph & 4 & \texttt{/coach/analyze} \\
Exam & Pipeline & Gemini Flash/Pro & StateGraph & 10 & \texttt{/exam/analyze} \\
TP Création & Pipeline & Gemini Flash & StateGraph & 5 & \texttt{/practical-work/\{id\}} \\
TP Correction & Pipeline & Gemini Flash & StateGraph & 2 & \texttt{/submissions/\{id\}/grade} \\
Feedback & Génératif & Gemini Flash & Appel direct & 1 & \texttt{/feedback/generate/\{id\}} \\
RAG/Chat & Retrieval & Gemini + ChromaDB & CAG$\rightarrow$RAG & 2 & \texttt{/ai/chat/\{doc\_id\}} \\
Program Extract & Pipeline & Gemini Flash & Multi-agents & 5+ & \texttt{/programs/\{id\}/process} \\
\bottomrule
\end{longtable}

% ============================================================================
\section{Stratégie de sélection des modèles}

\begin{table}[h!]
\centering
\begin{tabular}{@{}p{4cm}p{5cm}p{5cm}@{}}
\toprule
\textbf{Modèle} & \textbf{Cas d'usage} & \textbf{Caractéristiques} \\
\midrule
Gemini 2.5-Flash & Chat, quiz, feedback (défaut) & Rapide, <3s, économique \\
Gemini 2.5-Pro & Examens, analyse approfondie & Raisonnement profond \\
TunBERT & Dialecte tunisien & 768 dim, classification d'intentions \\
ChromaDB embeddings & RAG documentaire & Recherche sémantique vectorielle \\
\bottomrule
\end{tabular}
\caption{Stratégie de sélection des modèles LLM}
\end{table}

\end{document}
